\section{Prerequisites}
Before starting, ensure that the following hardware and software components are available:

\subsection{Hardware}
\begin{itemize}
    \item \textbf{ChipWhisperer Lite (CW1173)}: The consolidated fault injection hardware and target platform.
    \item \textbf{STM32F303 Target}: The microcontroller (integrated on the Lite board) running the Kyber implementation.
\end{itemize}

\subsection{Software}
\begin{itemize}
    \item \textbf{ARM Toolchain}: \texttt{arm-none-eabi-gcc} for compiling the firmware.
    \item \textbf{Python 3}: For running attack scripts.
    \item \textbf{ChipWhisperer Software}: Python library for controlling the hardware.
\end{itemize}

\section{Installation Steps}

\subsection{1. Clone the Repository}
The base library used is \texttt{pqm4}, which contains post-quantum cryptography schemes for the Cortex-M4.

\begin{lstlisting}[language=bash]
git clone https://github.com/mupq/pqm4.git
cd pqm4
git submodule update --init --recursive
\end{lstlisting}

\subsection{2. Install Python Dependencies}
Install the required Python packages, including the ChipWhisperer library.

\begin{lstlisting}[language=bash]
pip install -r requirements.txt
\end{lstlisting}

\subsection{3. Compile the Firmware}
We created a specific test firmware \texttt{test\_glitch} to facilitate the attack. This firmware generates a keypair, creates an invalid ciphertext, and then enters an infinite loop of decapsulation attempts.

To compile the firmware for the STM32F3 platform:

\begin{lstlisting}[language=bash]
make PLATFORM=cw308t-stm32f3 test-glitch
\end{lstlisting}

This command generates the binary file \texttt{bin/test\_glitch.hex}, which can be flashed onto the target device.
